%! suppress = EscapeUnderscore
In the previous chapter we showed how given the probability distribution of a random variable, we know everything
about it, and we can compute probabilities, expected values, and many more. In a way this is the job of descriptive
statistics. In this chapter we will show how we can draw conclusions for the population, when the population is not
accessible. This is the job of inferential statistics. \v

Formally, statistical inference is the process of using data analysis to infer properties of an underlying
distribution of probability. Inferential statistical analysis infers properties of a population, for example by
testing hypotheses and deriving estimates. It is assumed that the observed data set is sampled from a larger
population. \v

Inferential statistics can be contrasted with descriptive statistics. Descriptive statistics is solely concerned with
properties of the observed data, and it does not rest on the assumption that the data come from a larger population.

\section{Population VS Sample}

Let's begin with some basic definitions.
\bd[Population]
A \textbf{population} is a set of similar items or events which is of interest for some question or experiment. A
statistical population can be a group of existing objects or a hypothetical and potentially infinite group of objects
conceived as a generalization from experience.
\ed

Up to this point, technically we have been talking for populations. A population is the general category under
examination. An r.v represents the population and the corresponding probability distribution tells us about the
behaviour of the r.v and subsequently the behaviour of the population.

\bd[Parameter]
A \textbf{parameter} is a characteristic of a population.
\ed

Some examples of parameters are the expected value (or mean), the variance and the standard deviation of the
population. \v

More often than not, the population is not available to us either because gathering data for all the population is
very expensive or in most of the cases because it is impossible.

\be
For example, if we assume that our population is men's weights, and we want to know about the mean parameter (i.e.\ the
mean weight of all men), weighting all men around the globe at the same time is impossible. Usually, we end up with a
very small portion of the population called a ``sample''. (To not be confused with ``sample space'')
\ee

\bd[Sample]
A \textbf{sample} is a subset of a population. The elements of a sample are known as sample points, sampling units or
observations.
\ed

\vspace{-15pt}

\fig{population_and_sample}{0.25}

\vspace{-15pt}

Similarly to the definition of a parameter:

\bd[Statistic]
A \textbf{statistic} is a characteristic of a sample.
\ed

It follows from the definition of the sample that the latter is collected out of a population.

\be
In our previous example, one potential sample could be the weights of 1000 men.
\ee

\subsection{Sampling Techniques}

Since samples are used to draw conclusion for the population one has to be extremely careful while collecting samples
in order for the sample to be a good representative of the population. More specifically we must make sure that
members of samples are randomly selected (each member of a population has an equal chance to be selected) and samples
themselves are randomly selected (each sample of a population has an equal chance of being selected). Here are some
definitions based on the way of collecting a sample. \v

First of all there are two big categories: ``non-probability'' and ``probability-based'' sampling.

\bd[Non-Probability Sampling]
\textbf{Non-Probability sampling} is when the selection of data isn't based on any probability criteria.
\ed

\bd[Probability-Based Sampling]
\textbf{Probability-based sampling} is when the selection of data is based on specific probability criteria.
\ed

Starting with non-probability sampling here are a few non-probability sampling methods.

\bd[Convenience Sampling]
\textbf{Convenience sampling} is a type of non-probability sampling that involves the sample being drawn from that
part of the population that is close to hand.
\ed

\bd[Snowball Sampling]
\textbf{Snowball sampling} is a type of non-probability sampling where future samples are selected based on existing
samples.
\ed

\bd[Judgment Sampling]
\textbf{Judgment sampling} is a type of non-probability sampling where experts decide what samples to include.
\ed

\bd[Quota Sampling]
\textbf{Quota sampling} is a type of non-probability sampling where you select samples based on quotas for certain
slices of data without any randomization.
\ed

In general, the samples selected by non-probability criteria are not representative of the real-world data and therefore
are riddled with selection biases. For this reason people use almost exclusively probability-based sampling techniques
which we will cover now.

\bd[Random Sampling]
\textbf{Random sampling} is a procedure for sampling from a population in which the selection of a sample unit is
based on chance and every element of the population has a known, non-zero probability of being selected
\ed

\bd[Stratified Sampling]
\textbf{Stratified sampling} is a probability sampling technique that divides the entire population into different
subgroups or strata, and then randomly selects the final subjects proportionally from the different strata.
\ed

\bd[Weighted Sampling]
\textbf{Weighted sampling} is a probability sampling technique where each sample is given a weight, which determines the
probability of it being selected.
\ed

\bd[Reservoir Sampling]
\textbf{Reservoir sampling} is a family of randomized algorithms for choosing a simple random sample, without
replacement, of $k$ items from a population of unknown size $n$ in a single pass over the items.
\ed

\bd[Importance Sampling]
\textbf{Importance sampling} is a Monte Carlo method for evaluating properties of a particular distribution, while only
having samples generated from a different distribution than the distribution of interest.
\ed

In mathematical terms, given a probability distribution $P$, a random sample of length $n$ is a set of realizations
of $n$ independent, identically distributed random variables (i.i.d r.v's) with distribution $P$. A sample concretely
represents the results of $n$ experiments in which the same quantity is measured.

\be
In our example, if we want to estimate the mean weight of members of a particular population, we measure the heights of
$n$ individuals.
\ee

Each measurement is drawn from the probability distribution $P$ characterizing the population, so each measured
weight $x_{i}$ is the realization of an r.v $X$ with distribution $P$. Hence, mathematically a sample can be
represented as $\{ x_{1}, x_{2}, \ldots, x_{n} \}$. Given this representation we can define some characteristics
(statistics) of a sample.

\bd[Sample Size]
Given a sample of realizations of of $n$ i.i.d r.v's $\{ x_{1}, x_{2}, \ldots, x_{n} \}$, we call \textbf{sample
size} the direct count of the number of samples measured or observations being made $n$.
\ed

\bd[Sample Mean]
Given a sample of realizations of $n$ i.i.d r.v's $\{ x_{1}, x_{2}, \ldots, x_{n} \}$, we define the \textbf{sample
mean} or average $\bar{x}$ of the sample as the quantity:
\bse
\bar{x} = \frac{1}{n} \sum_{i=1}^{n} x_{i}
\ese
\ed

\bd[Sample Variance]
Given a sample of realizations of $n$ i.i.d r.v's $\{ x_{1}, x_{2}, \ldots, x_{n} \}$, we define the \textbf{sample
variance} $s^2$ of the sample as the quantity:
\bse
s^2 = \frac{1}{n} \sum_{i=1}^{n} (x_{i} - \bar{x})^2
\ese
\ed

\bd[Sample Standard Deviation]
Given a sample of realizations of $n$ i.i.d r.v's $\{ x_{1}, x_{2}, \ldots, x_{n} \}$, we define the \textbf{sample
standard deviation} $s$ of the sample as the quantity:
\bse
s = \sqrt{s^2} = \sqrt{\frac{1}{n} \sum_{i=1}^{n} (x_{i} - \bar{x})^2}
\ese
\ed

\bt[Law Of Large Numbers]
Given a random variable $X$ that follows a probability distribution $P$ with mean $\mu$ and a collection of $n$
realizations of $X$ $\{ x_{1}, x_{2}, \ldots, x_{n} \}$ with an average $\bar{x}_{n}$ then:
\bse
\bar{x}_{n} \to \mu \:\:\: for \:\:\: n \to \infty
\ese
\et

The law of large numbers is a theorem that describes the result of performing the same experiment a large number of
times. According to the law, the average of the results obtained from a large number of trials should be close to the
actual expected value of the population, and will tend to become closer to the expected value as more trials are
performed.The law of large numbers is important because it guarantees stable long-term results for the averages of
some random events.

\fig{lawoflargenumbers}{0.24}

The law of large numbers can find an application in statistics since, the ``random variable $X$ that follows a
probability distribution $P$ with mean $\mu$'' can be translated to a population, and the ``collection of $n$
realizations of $X$ $\{ x_{1}, x_{2}, \ldots, x_{n} \}$ with an average $\bar{x}_{n}$'' can be translated to a sample of
size $n$ drawn out of the population. Then the law of large numbers simply states that as the size of the sample
increases the sample mean tends to the actual population mean.

\be
Going back to our example, this means that as the number of men in our sample increases their average weight tends to
the actual mean weight of the whole population.
\ee

Hence, when we want to study a parameter of a population, (given that the population is huge and inaccessible in its
totality) we choose a random sample out of the population, following the techniques we have already introduced, and
we study the statistics on it. Then we generalize the statistic from the sample to the parameter of the population.
inevitably a question arises: ``how sure can we be that the statistics is close to the parameter''? \v

Practically one could say that the solution in order to obtain the correct parameter of a population out of a sample
is to increase the sample size as much as possible., so the law of large numbers will apply, and we get the correct
parameter. However, this is not true since the sample size follows the so-called ``law of diminishing returns, i.e.\
there is a point when increasing the sample size even more does not offer any statistical significance. \v

The solution to this problem is the so called ``central limit theorem'' and something called ``sampling distribution''.

\section{Central Limit Theorem \& Sampling Distribution}

\bt[Central Limit Theorem (CLT)]
Let $\{X_{1},\ldots,X_{n}\}$ be a sequence of i.i.d r.v's with $E[X_{i}]=\mu$ and $Var[X_{i}] = \sigma^{2} < \infty$.
Then as $n$ approaches infinity, the random variables $ \sqrt{n} ({\bar {X}}_{n} - \mu)$ converge in distribution
to a normal distribution $N(0,\sigma ^{2})$.
\et

CLT establishes that when independent random variables are added, their properly normalized sum tends toward a normal
distribution (informally a bell curve) even if the original variables themselves are not normally distributed. The
theorem is a key concept in probability theory because it implies that probabilistic and statistical methods that
work for normal distributions can be applicable to many problems involving other types of distributions. \v

For example, suppose that a sample is obtained containing many observations, each observation being randomly
generated in a way that does not depend on the values of the other observations, and that the arithmetic mean of the
observed values is computed. If this procedure is performed many times, the central limit theorem says that the
probability distribution of the average will closely approximate a normal distribution.

\be
A simple example of this is that if one flips a coin many times, the probability of getting a given number of heads
will approach a normal distribution, with the mean equal to half the total number of flips. At the limit of an
infinite number of flips, it will equal a normal distribution.
\ee

Now we will make use of CLT in order to develop statistical inference. Let's begin with some needed definitions.

\bd [Sampling Distribution]
A \textbf{sampling distribution} is the probability distribution of a given random sample based statistic.
\ed

If an arbitrarily large number of samples, each involving multiple observations, were separately used in order to
compute one value of a statistic (such as, for example, the sample mean or sample variance) for each sample, then
the sampling distribution is the probability distribution of the values that the statistic takes on. In many
contexts, only one sample is observed, but the sampling distribution can be found theoretically. \v

Although everything can be applied for any statistics, for now we will focus on the population mean (and
subsequently to the sample mean). Hence, given the sampling distribution of sample means (aka a probability
distribution of sample means) we know from CLT that it will follow a normal distribution independently of the
probability distribution of the population. Moreover, from the law of large numbers:
\bse
E[\bar{X}_{n}] = \mu, \quad \text{as } n \to \infty
\ese

\bd [Standard Error Of The Mean]
The \textbf{standard error of the mean}, or simply standard error, is defined as:
\bse
\sigma_{\bar {x}} = \frac {\sigma }{\sqrt {n}}
\ese

where $\sigma$ is the standard deviation of the population and $n$ is the size of the sample.
\ed

Observe that for $n \to \infty \Rightarrow \sigma_{\bar {x}} \to 0$. \v

Summing up according to CLT the sampling distribution of the mean of a population follows a normal distribution with
mean $E[\bar{X}_{n}]$ and and variance $\sigma_{\bar {x}}$.

\section{Confidence Intervals}

Since according to CLT the sampling distribution of the mean follows a normal distribution, we can use everything we
know about normal distributions. For example we know that 95\% of the values of a normal distribution lie within
1.96 standard deviations around the mean. For the sampling distribution of the mean this reads: 95\% of all the
samples that one could pick out of the population will have a mean $\bar{x}$ inside the interval:
\bse
[E[\bar{x}] - 1.96 \sigma_{\bar{x}}, \: E[\bar{x}] + 1.96 \sigma_{\bar{x}}]
\ese

\fig{confidenceinterval}{0.6}

The rest of the area (the uncovered one) is called ``significance level'' and it is denoted by $\alpha$. For
example, for a 95\% confidence interval we have $\alpha = 5\%$.

\bd[Significance Level]
The \textbf{significance level} $\alpha$ is the percentage of the sampling distribution of a statistic that is not
covered by the corresponding confidence interval.
\ed

In other words, for 95\% confidence intervals, the 95\% of the samples of a population will have mean value:
\bse
\bar{x} - 1.96 \frac{\sigma}{\sqrt{n}} \leq \mu \leq \bar{x} + 1.96 \frac{\sigma}{\sqrt{n}}
\ese

This leads us to the following definition.

\bd[Confidence Interval]
Given a sample of size $n$ with sample mean $\bar{x}$, drawn out of a population with standard deviation $\sigma$,
the \textbf{confidence interval} of $(1-\alpha)$ level of confidence for the population mean $\mu$ is defined as:
\bse
\bar{x} \pm z_{\frac{\alpha}{2}} \frac{\sigma}{\sqrt{n}}
\ese

where $z_{\frac{\alpha}{2}}$ is the $z$-score that leaves an area of $\frac{\alpha}{2}$ on each tail of the standard
normal distribution.
\ed

\textbf{Be careful:} A confidence interval is \emph{not} the probability that the population mean lies within the
interval! The population mean $\mu$ is a fixed (yet unknown) number, and the interval is the random object: if we
repeat the sampling many times, $(1-\alpha)$ of the confidence intervals we construct will contain $\mu$. \v

Observe that in the definition of the confidence interval we used the standard deviation $\sigma$ of the population,
which more often than not is unknown to us. In this case we use the standard deviation $s$ of the sample instead,
but then the sampling distribution of the mean is no longer a normal distribution. It follows the so-called
``Student's t-distribution'' with $n-1$ degrees of freedom, which we introduced in the previous chapter:
\bse
f_{X}(x) = \frac{\Gamma \Big( \frac{n+1}{2} \Big) }{\sqrt{n \pi}
\Gamma \Big( \frac{n}{2} \Big) } \Big ( 1 + \frac{x^2}{n} \Big)^{- \frac{n + 1}{2}}
\ese

In general the t-distribution is shorter in the middle and fatter in the tails compared to the normal distribution,
i.e.\ it gives a greater chance to extreme values, which compensates for the extra uncertainty of not knowing
$\sigma$. As $n \to \infty$ the t-distribution tends to the normal distribution. \v

Hence, when we have a t-distribution, for a given level of confidence, the confidence intervals will be given by:
\bse
\bar{x} \pm t_{cr} \frac{s}{\sqrt{n}}
\ese

where $s$ is the standard deviation of the sample, and $t_{cr}$ is a specific value (we can find it in t-tables)
which is specified by the level of confidence (or $\alpha$) and the degrees of freedom, i.e.\ the sample size minus
one.

\subsection{Margin Of Error \& Sample Size Determination}

The quantity that we add and subtract to the sample mean in order to form a confidence interval has a name of its
own.

\bd[Margin Of Error]
The \textbf{margin of error} $E$ of a confidence interval is the quantity:
\bse
E = z_{\frac{\alpha}{2}} \frac{\sigma}{\sqrt{n}}
\ese

so that the confidence interval reads $\bar{x} \pm E$. (Same for t, with $t_{cr}$ and $s$.)
\ed

Observe that the margin of error depends only on the sample size $n$, the level of confidence $\alpha$, and the
standard deviation of the population for $z$ (or of the sample for $t$). Since we cannot fix the standard
deviation, and usually the $\alpha$ is fixed to either $0.05$ or $0.01$, the only thing we can vary is the sample
size $n$. Solving the definition of the margin of error for $n$ we get:
\bse
n = \frac{z_{\frac{\alpha}{2}}^{2} \: \sigma^{2}}{E^{2}}
\ese

Hence, by specifying $z_{\frac{\alpha}{2}}$, $\sigma$ and $E$ we can come up with the suitable sample size for our
experiment.

\be
For 95\% confidence (i.e.\ $z_{\frac{\alpha}{2}} = 1.96$), given a fixed error $E$ that we don't care if we're
wrong, and a specific standard deviation $\sigma$, we need a sample with sample size $n = (1.96)^{2} \sigma^{2} /
E^{2}$.
\ee

Of course $n \to \infty$ as $E \to 0$, i.e.\ the more precision we require, the bigger the sample we have to collect.

\subsection{Confidence Interval For The Variance}

Everything we did so far concerns the population mean $\mu$, but the exact same logic can be applied to the
population variance $\sigma^{2}$ as well. The notation that we will use for the variances is the same that we have
been using from the beginning of the chapter:
\bse
\sigma^{2} = \frac{\sum_{i} (x_{i} - \mu)^2}{N} \quad \text{(population)} \qquad \qquad
s^{2} = \frac{\sum_{i} (x_{i} - \bar{x})^2}{n-1} \quad \text{(sample)}
\ese

and for the standard deviations $\sigma = \sqrt{\sigma^{2}}$ (population) and $s = \sqrt{s^{2}}$ (sample). \v

As we discussed, the distribution of the means taken by multiple samples follows a normal distribution
$N \left( \mu, \sfrac{\sigma}{\sqrt{n}} \right)$. Similarly, in the sampling distribution of the variance, the
corresponding statistic for a sample of sample size $n$ with sample variance $s^{2}$ is the quantity:
\bse
\chi^{2} = \frac{(n-1) s^{2}}{\sigma^{2}}
\ese

which follows a chi-squared distribution with $n-1$ degrees of freedom, which we also introduced in the previous
chapter. Of course, since variation cannot be negative, it makes sense for it to follow a distribution with a zero
PDF for $x < 0$. \v

Observe that, contrary to the normal distribution where everything was symmetrical around the mean, the chi-squared
distribution is not symmetrical. So when we define an $\alpha$, it does not split the distribution in equal parts,
and instead of one $z_{cr}$ we now have to find two different values $\chi^{2}_{1 - \frac{\alpha}{2}}$ and
$\chi^{2}_{\frac{\alpha}{2}}$ from a $\chi^{2}$-table, for the given degrees of freedom $n - 1$. (As a convention,
in $\chi^{2}$ the cumulative probability 1 starts from the left and 0 from the right, opposite of the normal
distribution which was 0 on the left and 1 on the right.) \v

After doing that, we know that $(1-\alpha)$ of all samples will have a $\chi^{2}$ between the two limits that we
found:
\bi{rCl}
\chi^{2}_{1 - \frac{\alpha}{2}} \: \leq \: \frac{(n-1) s^{2}}{\sigma^{2}} \: \leq \: \chi^{2}_{\frac{\alpha}{2}}
& \eqv &
\frac{1}{\chi^{2}_{1 - \frac{\alpha}{2}}} \: \geq \: \frac{\sigma^{2}}{(n-1) s^{2}} \: \geq \:
\frac{1}{\chi^{2}_{\frac{\alpha}{2}}} \\
& \eqv & \frac{(n-1) s^{2}}{\chi^{2}_{\frac{\alpha}{2}}} \: \leq \: \sigma^{2} \: \leq \:
\frac{(n-1) s^{2}}{\chi^{2}_{1 - \frac{\alpha}{2}}}
\ei

So, we estimated the variance! Exactly the same thing as we did with means. (For standard deviations we just take
the square root of the inequality.) \v

Observe that as we change the sample size $n$ we change the degrees of freedom, hence the
$\chi^{2}_{1 - \frac{\alpha}{2}}$ and $\chi^{2}_{\frac{\alpha}{2}}$. In general as $n \to \infty$ the confidence
interval of $\sigma^{2}$ narrows. The idea is that as $n \to \infty$ the interval converges to a specific value
which is the ``real'' value of $\sigma^{2}$. However, the law of diminishing returns applies also here: after a
point, further increase to $n$ does not offer anything.

\section{Hypothesis Testing}

Confidence intervals are one of the two pillars of statistical inference: they estimate a population parameter out
of a sample. The second pillar answers a slightly different question: given a claim (a ``hypothesis'') about a
population parameter, does a sample support the claim or not? This process is called ``hypothesis testing'' and, as
we will see, it heavily relies on the sampling distributions and the confidence intervals we have developed so far.

\subsection{Null \& Alternative Hypotheses}

Let's begin with the two fundamental definitions.

\bd[Null Hypothesis]
The \textbf{null hypothesis} $H_{0}$ is the default position of a hypothesis test that there is nothing new or
different, hence what it's given (or assumed), i.e.\ the status quo is correct, and it is assumed true. It always
contains an equality.
\ed

\bd[Alternative Hypothesis]
The \textbf{alternative hypothesis} $H_{a}$ is the contrary to the null hypothesis, hence the other option when the
null hypothesis is rejected. Usually it's unknown, it's not given, and it's a claim we make opposite to the null
hypothesis. It never contains an equality, it just rejects the equality of the null hypothesis, hence it proves the
opposite inequality.
\ed

Since the null hypothesis is an assumed fact, we accept it. A statistical hypothesis test will either
\textbf{reject} or \textbf{fail to reject} the null hypothesis. If we reject the null hypothesis it means that the
alternative hypothesis is true. However, if we fail to reject the null hypothesis, it does not mean that we have
proven that it's true, because we only \emph{assumed} it was true! \v

\textbf{Be careful:} We can never prove that the null hypothesis is true. We can only do that for the alternative
hypothesis! \v

Using this idea it's easy to develop the null and alternative hypothesis for any experiment.

\subsection{Type I \& Type II Errors}

Since hypothesis testing is based on samples, and samples carry randomness, there are two distinct ways in which
our conclusion can be wrong.

\bd[Type I Error]
A \textbf{Type I error} is the rejection of the null hypothesis when it should have been accepted. In other words:
the incorrect rejection of a true null hypothesis.
\ed

\bd[Type II Error]
A \textbf{Type II error} is the failure to reject the assumption (null hypothesis) when it should have been
rejected. In other words: the incorrect acceptance of a false null hypothesis.
\ed

The most usual way to do a Type II error is by randomly selecting a sample that approves the null hypothesis.
Schematically:

\v

\btab[H]
\centering
\btb{c|c||c|c}
& & \multicolumn{2}{c}{Conclusion} \\
& & Accept $H_{0}$ & Reject $H_{0}$ \\
\hline
\rule{0pt}{12pt} \multirow{2}{*}{Actual Condition} & $H_{0}$ is true & \checkmark & Type I error \\
& $H_{a}$ is true & Type II error & \checkmark
\etb
\etab

\v

In real life Type II errors are much more important than Type I errors, since failing to detect a real change is
usually more costly than raising a false alarm. \v

We can use hypothesis testing in combination with the sampling distribution, in order to find if there is a
statistically significant change in a statistic or not. This is called a ``z-test'' or a ``t-test'' depending on
the distribution.

\subsection{The z-Test \& t-Test}

Here we are focusing on the process of a z-test, but the same applies for a t-test if we change all $z$-values to
$t$-values and all population standard deviations to sample standard deviations. \v

Given an assumed population mean $\mu_{0}$, we form the null hypothesis to be:
\bse
H_{0}: \mu = \mu_{0}
\ese

i.e.\ the actual population mean is the given one. Of course the alternative hypothesis is:
\bse
H_{a}: \mu \neq \mu_{0}
\ese

We develop the test given an $\alpha$, i.e.\ a ``significance level'' as we defined it for confidence intervals,
which now obtains one more interpretation: that of a ``Type I error rate''. It states how tolerant we are to make a
Type I error. In other words, what percentage of total samples should we accept to produce a Type I error.

\be
For $\alpha = 5\%$, I am okay to make a Type I error 5 out of 100 times.
\ee

We call the corresponding $z$ value for a given $\alpha$ the ``critical value''.

\bd[Critical Value]
The \textbf{critical value} $z_{cr}$ (or $t_{cr}$) is the value of the test statistic that corresponds to the
chosen significance level $\alpha$, i.e.\ the boundary beyond which the null hypothesis is rejected. The critical
value is determined by $\alpha$ and by whether we are using the z or the t distribution.
\ed

Let's see this in a graph!

\fig{criticalregion}{0.6}

We know that $(1-\alpha)$ of all random samples will have a mean inside the region $[\mu_{0} - z_{cr} \sigma_{\bar
{x}}, \: \mu_{0} + z_{cr} \sigma_{\bar{x}}]$ defined by the critical values. The rest $\alpha$ of the samples will
have a mean outside this region, but they will still be under the same distribution! \v

Now, we assume that the population mean has changed to some $\mu \neq \mu_{0}$ and we want to test that, so we form
the null and the alternative hypotheses as above. We gather a sample of size $n$, we compute its mean $\bar{x}$,
and we want to check if $\bar{x}$ is inside the non-rejection area:
\bi{rCl}
\mu_{0} - z_{cr} \sigma_{\bar{x}} \: \leq \: \bar{x} \: \leq \: \mu_{0} + z_{cr} \sigma_{\bar{x}}
& \eqv &
- z_{cr} \: \leq \: \frac{\bar{x} - \mu_{0}}{\sigma_{\bar{x}}} \: \leq \: z_{cr}
\ei

where the middle quantity is nothing but the $z$-score of $\bar{x}$:
\bse
z = \frac{\bar{x} - \mu_{0}}{\sfrac{\sigma}{\sqrt{n}}} \qquad \Big( \text{or the t-score of it:} \quad
t = \frac{\bar{x} - \mu_{0}}{\sfrac{s}{\sqrt{n}}} \Big)
\ese

Note that we use $\sfrac{\sigma}{\sqrt{n}}$ (or $\sfrac{s}{\sqrt{n}}$) because the sample has a sample size $n$,
hence $\bar{x}$ lives in the sampling distribution of the mean. \v

Even if the population mean is no longer $\mu_{0}$, the $z$-score of our sample might be between $-z_{cr}$ and
$z_{cr}$, so it's still part of the previous distribution and $H_{0}$ cannot be rejected. However, if
$z < -z_{cr}$ or $z > z_{cr}$, then there are two possible cases:

\bit
\item $z$ is part of the $\alpha$ of the random samples that have a mean outside $[-z_{cr}, z_{cr}]$, hence the
distribution has not changed, and we are performing a Type I error.
\item $z$ is part of the $(1 - \alpha)$ of the random samples of a new distribution, so the distribution has
changed!
\eit

Hence the decision rule is:
\bse
z < -z_{cr} \:\: \text{or} \:\: z > z_{cr} \imp \text{Reject } H_{0} \qquad \qquad
-z_{cr} \leq z \leq z_{cr} \imp \text{Fail to reject } H_{0}
\ese

and when we reject $H_{0}$, we have an $\alpha$ probability of being wrong, i.e.\ of performing a Type I error. \v

It makes sense that the smaller the $\alpha$, the bigger the critical value $z_{cr}$, since the more sure we want
to be, the broader the range of possible means for our samples. Hence as $\alpha \to 0$, while it's ``difficult''
to find a $z$ out of statistical significance ranges, once we find it, it's very unlikely to have a Type I error.
This is why $\alpha$ is the probability of having a Type I error! \v

So the smaller the $\alpha$, the more ``difficult'' to find a $z$ outside the critical range (i.e.\ a statistically
significant $z$), but once we find it, the less the probability of having committed a Type I error when we reject
the initial distribution. However, as we decrease $\alpha$, the probability of having a Type I error decreases,
but the probability of a Type II error increases, because now we are prone to pick a $z$ belonging to a
``further-up'' distribution, inside the critical range.

\fig{typeoneerrortradeoff}{0.6}

As we can see in the graph, a point right below the critical value is inside the non-rejection region of the null
distribution, but it may very well belong to the shifted distribution: failing to reject it is exactly a Type II
error.

\subsection{The p-Value}

An alternative way to decide for rejecting $H_{0}$ or not is using the ``p-value''. (Exactly the same way!)

\bd[p-Value]
We define as \textbf{p-value} the probability of finding the observed (or more extreme) results when the null
hypothesis is true.
\ed

In simple words, once we find the $z$ of the sample we calculate the ``remaining area'':

\fig{pvalue}{0.6}

Hence for the $z$ we calculated, we find through a z-table to which probability it corresponds. This way we find
the p-value! It's intuitive that in order to reject $H_{0}$ it must be:
\bse
p < \alpha \qquad (\text{which translates to} \:\: z < -z_{cr} \:\: \text{or} \:\: z > z_{cr})
\ese

Same result, different way of expressing it! (Same for t-tests. Just find the $\alpha$ of degrees of freedom $n-1$
and the $t$-value we calculated. If $p < \alpha$ then we reject $H_{0}$, otherwise we fail to reject it.)

\subsection{Hypothesis Testing For Two Populations}

So far we have been testing a single population against an assumed value $\mu_{0}$. However, very often we want to
compare two populations with each other. There are two distinct cases: the populations can be independent, or
dependent (matched). \v

Starting with the independent case, we are given two independent populations with mean values $\mu_{1}, \mu_{2}$
and standard deviations $\sigma_{1}, \sigma_{2}$. By ``independent'', we mean that measurements on one population
do not affect the other. We create our hypotheses as:
\bse
H_{0}: \mu_{1} - \mu_{2} = D_{0} \qquad \qquad H_{a}: \mu_{1} - \mu_{2} \neq D_{0}
\ese

where $D_{0}$ is the assumed difference of the means. (For example $D_{0} = 0$ tests whether the two populations
have the same mean.) The sample differences of the means $\bar{x}_{1} - \bar{x}_{2}$ satisfy a normal distribution
centered at $D_{0}$ with standard deviation:
\bse
\sigma^{*} = \sqrt{\frac{\sigma_{1}^{2}}{n_{1}} + \frac{\sigma_{2}^{2}}{n_{2}}}
\ese

where $n_{1}, n_{2}$ are the two sample sizes. Hence the corresponding $z$-score of the observed difference is:
\bse
z = \frac{(\bar{x}_{1} - \bar{x}_{2}) - D_{0}}{\sqrt{\frac{\sigma_{1}^{2}}{n_{1}} + \frac{\sigma_{2}^{2}}{n_{2}}}}
\ese

We pick an $\alpha$, and by choosing $\alpha$ we have chosen a $z_{cr}$. As before:
\bse
z < -z_{cr} \:\: \text{or} \:\: z > z_{cr} \imp \text{Reject } H_{0} \qquad \qquad
-z_{cr} \leq z \leq z_{cr} \imp \text{Fail to reject } H_{0}
\ese

and we have an $\alpha$ probability of being wrong and performing a Type I error. As before, the ``remaining''
uncovered area defined by $z$ is called the p-value, and it should be $p < \alpha$ for us to reject $H_{0}$. \v

In the case that we don't know $\sigma_{1}, \sigma_{2}$, we simply estimate them by using the sample standard
deviations $s_{1}$ and $s_{2}$ of the samples. Then we perform a t-test:
\bse
t = \frac{(\bar{x}_{1} - \bar{x}_{2}) - D_{0}}{\sqrt{\frac{s_{1}^{2}}{n_{1}} + \frac{s_{2}^{2}}{n_{2}}}}
\ese

and finally by fixing $\alpha$ we obtain a $t_{cr}$ and we repeat the same steps. \v

We just reviewed a case where the two populations were independent. However, we can develop a matched t-test for
two dependent populations. Usually this is the case when we have a population with a ``before'' and an ``after''
condition, and we want to study if the mean value of the two situations has changed statistically significantly.

\fig{matchedttest}{0.6}

Given the ``before'' measurements $(x_{1}, x_{2}, \ldots, x_{n})$ and the ``after'' measurements
$(x_{1}', x_{2}', \ldots, x_{n}')$ of the same $n$ units, we define the differences $d_{i} = x_{i} - x_{i}'$ for
all $i$, and we calculate the mean difference:
\bse
\bar{d} = E[d_{i}]
\ese

Under the null hypothesis that nothing has changed, the distribution of $\bar{d}$ has mean value 0, hence:
\bse
H_{0}: \bar{d} = 0 \qquad \qquad H_{a}: \bar{d} \neq 0
\ese

Now we follow the same procedure of a t-test. We pick a level of confidence (e.g.\ $\alpha = 0.05$), and since we
have $n$ sample size the degrees of freedom are $n-1$. Out of this we obtain a $t_{cr}$. We already know $\bar{d}$
and we also calculate $s_{d}$, i.e.\ the standard deviation of the differences. We find the t-value of $\bar{d}$:
\bse
t = \frac{\bar{d} - 0}{\sfrac{s_{d}}{\sqrt{n}}} = \frac{\sqrt{n} \: \bar{d}}{s_{d}}
\ese

\textbf{Careful:} What we did, is that we assumed that the initial distribution of $\bar{d}$ (the one we compare
our sample right now) had mean value 0 (this is why we subtract 0), and since we're doing a t-test it had the
standard deviation of the given samples. \v

So:
\bse
t < -t_{cr} \:\: \text{or} \:\: t > t_{cr} \imp \text{Reject } H_{0} \qquad \qquad
-t_{cr} \leq t \leq t_{cr} \imp \text{Fail to reject } H_{0}
\ese

\subsection{Hypothesis Testing For The Variance}

Exactly as we did for confidence intervals, hypothesis testing can also be applied to variances instead of means.
We assume that the variance has changed, and we form the hypotheses:
\bse
H_{0}: \sigma^{2} = \sigma_{0}^{2} \qquad \qquad H_{a}: \sigma^{2} \neq \sigma_{0}^{2}
\ese

We pick a sample of sample size $n$, and we measure its variance $s^{2}$. We pick a level of confidence $\alpha$,
and by using degrees of freedom $n-1$, we obtain $\chi^{2}_{1 - \frac{\alpha}{2}}$ and
$\chi^{2}_{\frac{\alpha}{2}}$ from a $\chi^{2}$-table. We calculate the $\chi^{2}$ of our sample as:
\bse
\chi^{2} = \frac{(n-1) s^{2}}{\sigma_{0}^{2}}
\ese

and the decision rule is:
\bse
\chi^{2}_{1 - \frac{\alpha}{2}} \leq \chi^{2} \leq \chi^{2}_{\frac{\alpha}{2}} \imp \text{Fail to reject } H_{0}
\qquad \qquad
\chi^{2} < \chi^{2}_{1 - \frac{\alpha}{2}} \:\: \text{or} \:\: \chi^{2} > \chi^{2}_{\frac{\alpha}{2}}
\imp \text{Reject } H_{0}
\ese

For two samples $(n_{1}, s_{1}^{2})$ and $(n_{2}, s_{2}^{2})$ we define the F-ratio as:
\bse
F = \frac{s_{1}^{2}}{s_{2}^{2}}
\ese

which follows the so-called ``F-distribution'' with $n_{1} - 1$ and $n_{2} - 1$ degrees of freedom, and we repeat
the same steps: fix an $\alpha$, obtain the critical values from an F-table, and compare.

\subsection{The Chi-Squared Goodness Of Fit Test}

The chi-squared distribution provides us with one more very useful test, which checks how well an observed
distribution of categorical data fits a theoretically expected one. In order to do this we define:

\bit
\item $O_{i}$: the number of observations of type $i$
\item $E_{i}$: the number of observations of type $i$ expected theoretically
\item $i$: the types of observations
\eit

Given this we calculate the $\chi^{2}$ of our sample as:
\bse
\chi^{2} = \sum_{i} \frac{(O_{i} - E_{i})^{2}}{E_{i}}
\ese

As before, we have degrees of freedom equal to the number of types minus one, and by choosing a level of confidence
$\alpha$, we can obtain through a $\chi^{2}$-test-table the value of $\chi^{2}_{cr}$. This value $\chi^{2}_{cr}$ is
the maximum value of the $\chi^{2}$-value of variance that we can accept due to chance, being $(1 - \alpha)$
confident. Hence:
\bse
\chi^{2} \leq \chi^{2}_{cr} \imp \text{Variation due to chance} \qquad \qquad
\chi^{2} > \chi^{2}_{cr} \imp \text{Variation \emph{not} due to chance}
\ese

\section{Analysis Of Variance (ANOVA)}

Suppose now that we want to compare the means of more than two populations. One idea would be to test them in
pairs with the tools we already have:
\bse
H_{0}: \bar{x}_{1} = \bar{x}_{2}, \:\: \alpha = 0.05 \qquad
H_{0}: \bar{x}_{1} = \bar{x}_{3}, \:\: \alpha = 0.05 \qquad \ldots
\ese

However, following this procedure the alphas are compound together, hence in the end the final level of confidence
would be:
\bse
(0.95) \cdot (0.95) \cdots (0.95) = (0.95)^{\binom{N}{2}}
\ese

where $N$ is the number of populations. This number decreases significantly as $N$ grows.

\be
For $N = 3$: $(0.95)^{3} \approx 0.85$, which gives $\alpha = 0.15$ instead of $\alpha = 0.05$. As $N$ increases,
the probability of performing a Type I error increases as well.
\ee

In order to avoid this, we use another method for comparing means of more than 2 populations which is called
``Analysis Of Variance'' or ``ANOVA''! So, given $N$ means $\bar{x}_{1}, \bar{x}_{2}, \ldots, \bar{x}_{N}$ we want to
find out if they are all part of the same distribution or they belong to different distributions. In order to do so,
first we have to introduce two concepts: the variability among (or between) the sample means, and the variability
around (or within) the distributions.

\bd[Variance Between]
The \textbf{variance between} is the variability among the sample means, i.e.\ the distance of each sample mean
from the grand mean of all the samples.
\ed

\bd[Variance Within]
The \textbf{variance within} is the variability around (inside) each distribution. We measure the spread (or
width) of each sample separately, not in relation with something. This internal spread is why it's called
``within'' distributions.
\ed

\fig{anova}{0.6}

In ANOVA we are dealing with the ratio:
\bse
\frac{\text{Variance Between}}{\text{Variance Within}} = \frac{\text{Distance from grand mean}}{\text{Internal
spread}}
\ese

Of course:
\bse
\text{Total Variance} = \text{Variance Between} + \text{Variance Within}
\ese

The general idea of ANOVA is:
\bse
\frac{\text{Variance Between}}{\text{Variance Within}} \gg 1 \imp \text{Reject } H_{0} \qquad \qquad
\frac{\text{Variance Between}}{\text{Variance Within}} \lesssim 1 \imp \text{Fail to reject } H_{0}
\ese

That was the general concept of ANOVA. Now let's see the different ANOVAs in detail.

\subsection{One-Way ANOVA}

In ANOVA the hypotheses are always of the form:
\bse
H_{0}: \text{The means are coming from the same distribution (equal)}
\ese
\bse
H_{a}: \text{The means are coming from different distributions (unequal)}
\ese

In One-Way ANOVA we have a single ``factor'' which can take some values and creates different groups (or columns).

\be
For the year of the students: $1, 2, 3, \ldots, C$, each year is a group, and we study e.g.\ the grades of each
group.
\ee

For each of the $C$ groups we collect a sample of sample size $n$ and we find the mean value:
\bse
F_{1} \to \text{Sample}_{1} \to \bar{x}_{1} \qquad
F_{2} \to \text{Sample}_{2} \to \bar{x}_{2} \qquad \ldots \qquad
F_{C} \to \text{Sample}_{C} \to \bar{x}_{C}
\ese

Then we find the grand mean $\bar{\bar{x}}$ by grouping all the samples together and calculating the mean value. \v

Now we define three deviations:

\bit
\item Sum of Squares Total: $\ds \text{SST} = \sum_{\text{all data}} (x_{i} - \bar{\bar{x}})^{2}$ \:\: (Total)
\item Sum of Squares Column: $\ds \text{SSC} = n \sum_{j=1}^{C} (\bar{x}_{j} - \bar{\bar{x}})^{2}$ \:\: (Between)
\item Sum of Squares Error: $\ds \text{SSE} = \sum_{j=1}^{C} \sum_{i=1}^{n} (x_{ij} - \bar{x}_{j})^{2}$ \:\:
(Within)
\eit

so that:
\bse
\text{Total Variance} = \text{Variance Between} + \text{Variance Within} \qquad \text{or} \qquad
\text{SST} = \text{SSC} + \text{SSE}
\ese

Now, using SST, SSC and SSE, and by keeping in mind that we have $C$ samples (based on the factor) of $n$ sample
size each, we define:

\bit
\item Degrees of freedom of columns (or groups): $\text{dof}_{C} = C - 1$
\item Degrees of freedom of error (or observations): $\text{dof}_{E} = C(n - 1)$
\eit

Using the degrees of freedom we define:

\bit
\item Mean Sum of Squares Columns: $\ds \text{MSC} = \frac{\text{SSC}}{\text{dof}_{C}}$
\item Mean Sum of Squares Error: $\ds \text{MSE} = \frac{\text{SSE}}{\text{dof}_{E}}$
\eit

Finally we define:
\bse
F = \frac{\text{MSC}}{\text{MSE}} \qquad \Big( = \frac{\text{Variance Between}}{\text{Variance Within}} \Big)
\ese

which follows an F-distribution, and given a level of confidence $\alpha$ we obtain an $F_{cr}$ from an F-table.
Hence:
\bse
F > F_{cr} \imp \text{Reject } H_{0} \qquad \qquad F \leq F_{cr} \imp \text{Fail to reject } H_{0}
\ese

where rejecting $H_{0}$ means that the means are coming from different populations. \v

And this is One-Way ANOVA\@.

\subsection{Two-Way ANOVA}

In Two-Way ANOVA we have multiple samples that differ based on a factor, and each observation in each sample
differs based on another factor.

\be
Grades for a course of specific students in different years:
\bit
\item Factor 1: Years (columns). Splits into different samples (columns).
\item Factor 2: Student (rows). Each sample contains the same students (rows).
\eit
\ee

This specific Two-Way ANOVA is called ``Two-Way ANOVA without replications'' because each student appears in each
year only once. If a student could appear twice (e.g.\ students may retake the exam in the same year) then we'd
have ``Two-Way ANOVA with replications''. First let's check without replication.

\subsubsection{Two-Way ANOVA Without Replications}

As before, we can calculate each individual sample mean $\bar{x}_{1}, \bar{x}_{2}, \ldots, \bar{x}_{C}$, and the
grand mean $\bar{\bar{x}}$. But now we can also calculate the mean of each row, since they are split based on a
common factor (e.g.\ student):
\bse
\bar{x}_{R_{1}}, \: \bar{x}_{R_{2}}, \: \ldots, \: \bar{x}_{R_{B}}
\ese

So we have:

\bit
\item $\bar{x}_{1}, \bar{x}_{2}, \ldots, \bar{x}_{C}$: for groups (columns)
\item $\bar{x}_{R_{1}}, \bar{x}_{R_{2}}, \ldots, \bar{x}_{R_{B}}$: for entries (rows)
\item $\bar{\bar{x}}$: overall mean (grand)
\eit

In the previous case we didn't have the second one! As before we compute SST and SSC, but now we have an extra
deviation:

\bit
\item Sum of Squares Block: $\ds \text{SSB} = C \sum_{i=1}^{B} (\bar{x}_{R_{i}} - \bar{\bar{x}})^{2}$ \:\: (Rows)
\eit

So what we did is that we splitted the variation even more, in variation coming from the first factor (i.e.\ groups
or columns) SSC, variation coming from the second factor (i.e.\ blocks or rows) SSB, and the remaining error SSE\.
So now:
\bse
\text{SST} = \text{SSC} + \text{SSE} + \text{SSB}
\ese

Again, as before, we define the degrees of freedom, for the first two same as before:
\bse
\text{dof}_{C} = C - 1 \qquad \qquad \text{dof}_{B} = B - 1 \qquad \qquad \text{dof}_{E} = (C-1)(B-1)
\ese

and the mean sums of squares $\text{MSC} = \sfrac{\text{SSC}}{\text{dof}_{C}}$ and
$\text{MSE} = \sfrac{\text{SSE}}{\text{dof}_{E}}$, but now we also include the blocks:
\bse
\text{MSB} = \frac{\text{SSB}}{\text{dof}_{B}}
\ese

So again we define the F-ratio as:
\bse
F_{C} = \frac{\text{MSC}}{\text{MSE}}
\ese

but now we also have an F-ratio for blocks (rows):
\bse
F_{B} = \frac{\text{MSB}}{\text{MSE}}
\ese

So, for each of the two F's we can perform two single ``One-Way'' ANOVAs, but now we are also based on the
different factors!

\subsubsection{Two-Way ANOVA With Replications}

In a Two-Way ANOVA with replications, for each block (row) based on a factor, of each group (column) based on
another factor, we have multiple datapoints. In this case, for each block, we find the mean value of all
repetitions. By doing so, we end up to a ``Two-Way ANOVA without repetitions'', and everything is the same as
before! \v

But, since we now work with mean values, we also get some error from there. We call this error ``error from
interactions'' and using a programming language we can calculate it. We won't do it now! \v

This summarizes ANOVAs, and with it the classical statistics that we can do with the normal distribution and
similar distributions. In the next chapter we will revisit statistical inference in a more formal setting, and
develop the theory of parametric inference.